\documentclass[10pt,a4paper]{amsart}
\usepackage[margin=19mm]{geometry}
\usepackage{amsmath,amssymb,mathtools}
\usepackage[scr=rsfs,cal=euler]{mathalfa}
\usepackage{xfrac}
\usepackage[unicode,hidelinks,bookmarks=false]{hyperref}
\newcommand{\ie}{i.e.\ }
\newcommand{\cf}{cf.\ }
\newcommand{\resp}{respectively\ }
\newcommand{\Subsection}{\subsection}
\providecommand{\nobreakdash}{\nobreak\hbox{-}}
\setlength{\emergencystretch}{2em}
\multlinegap0pt
\usepackage{bm}
\usepackage{bbm}
\usepackage{dsfont}
\usepackage{xspace}
\usepackage{cancel}
\usepackage{mathtools}
\usepackage{amsthm}
\makeatletter
\@ifundefined{theorem}{\newtheorem{theorem}{Theorem}}{}
\@ifundefined{lemma}{\newtheorem{lemma}[theorem]{\normalfont \scshape Lemma}}{}
\makeatother
\makeatletter
\newcommand{\R}{\mathbb{R}}
\newcommand{\bbN}{\mathbb N}
\newcommand{\calO}{\mathcal O}
\newcommand{\calT}{\mathcal T}
\newcommand{\ds}{\textnormal{d}s}
\newcommand{\dt}{\textnormal{d}t}
\newcommand{\dv}{\textnormal{d}v}
\newcommand{\dx}{\textnormal{d}x}
\newcommand{\intr}{\int_{\R^d}}
\newcommand{\lt}{\left}
\newcommand{\p}{\partial}
\expandafter\ifx\csname proof\endcsname\relax
\renewenvironment{proof}[1][\proofname]{\noindent{\scshape #1.\quad}}{\qed}
\fi
\newcommand{\renorm}[3]{\lt[#1\rt]_{#2 , #3}}
\newcommand{\rt}{\right}
\newcommand{\tnd}{\textnormal{d}}
\newcommand{\ve}{\varepsilon}
\newcommand{\weakto}{\rightharpoonup}
\makeatother
\title[Global weak solutions to nonlinear kinetic Fokker--Planck eq]{Global weak solutions to nonlinear kinetic Fokker--Planck equations in bounded domains under physical initial data}
\author{Young-Pil Choi, Sihyun Song}
\date{Source: arXiv:2510.06656v1 (CC BY 4.0)}
\begin{document}
\maketitle

\noindent\emph{Source and licence.} Global weak solutions to nonlinear kinetic Fokker--Planck equations in bounded domains under physical initial data. Version arXiv:2510.06656v1 (CC BY 4.0), distributed under the Creative Commons Attribution 4.0 International licence, as recorded in the arXiv licence metadata for this exact version and shown on the abstract page https://arxiv.org/abs/2510.06656v1. Full rights notice: licenses/arxiv-2510.06656-CC-BY-4.0.txt.\par\smallskip

\noindent\textit{Editorial scope.} Counted unit: the Lemma whose source label is \texttt{eq: PASS weak form 1} in the article (source0.tex lines 1936--1941, source label \texttt{eq: PASS weak form 1}), delivered together with the article's own complete original proof of it (source0.tex lines 1942--2001, one \texttt{proof} environment of 60 lines). No mathematics is written, repaired or re-derived: statement and proof are the article's own text line for line. Required context retained uncounted: eq: reg (lines 1166--1172); lem: ve weak sol (lines 1432--1455); lem: subcubic (lines 1754--1756); eq: MAC convergence (lines 1827--1829); eq: V convergence (lines 1831--1833); eq: NU convergence (lines 1846--1848). Every definition, display and label of those lines is retained unchanged. Omitted from this package and not redistributed: 1--1165, 1173--1431, 1456--1753, 1757--1826, 1830--1830, 1834--1845, 1849--1935, 2002--3171. Nothing inside a retained excerpt is omitted or altered: every retained body is delivered exactly as the article's lines are written, so no omission receipt of any kind accompanies this package. Citations: the retained excerpts carry no citation command at all, so no bibliography entry of the article is transcribed and no reference is redistributed. Reference adaptation: none was needed and none was made; every reference inside the delivered unit resolves inside it, so no unresolved reference or citation is produced. Printed slips kept literal: no printed word, symbol, hypothesis or number of the counted statement or of its proof was altered. Layout and class: the article's own class is \texttt{amsart} with packages including amsmath, amsthm, amssymb, setspace, babel, mathrsfs, xy, paralist, graphicx, accents, enumerate, mathtools, csquotes, hyperref; the wrapper is the lane's pinned \texttt{amsart} editorial wrapper at 10pt on A4, which restates the theorem-like environments the retained text needs and adds the article's own macro definitions that the retained text needs (\texttt{\textbackslash R}, \texttt{\textbackslash bbN}, \texttt{\textbackslash calO}, \texttt{\textbackslash calT}, \texttt{\textbackslash ds}, \texttt{\textbackslash dt}, \texttt{\textbackslash dv}, \texttt{\textbackslash dx}, \texttt{\textbackslash intr}, \texttt{\textbackslash lt}, \texttt{\textbackslash p}, \texttt{\textbackslash renorm}, \texttt{\textbackslash rt}, \texttt{\textbackslash tnd}, \texttt{\textbackslash ve}, \texttt{\textbackslash weakto}), with extra wrapper packages (bm, bbm, dsfont, xspace, cancel, mathtools). The delivered numbering is the unit's own. Assets: no retained span contains a figure environment, a graphics-inclusion command, a PDF-page inclusion, a table or a TikZ picture; no image file is copied into this package, the package declares no asset dependency and no image file is redistributed with it. Licence: version arXiv:2510.06656v1 of this article is distributed under the Creative Commons Attribution 4.0 International licence, as recorded in the arXiv licence metadata for this exact version and shown on the abstract page https://arxiv.org/abs/2510.06656v1. A full rights notice is delivered at licenses/arxiv-2510.06656-CC-BY-4.0.txt. Characters: every retained excerpt is pure ASCII, so the delivered engine, XeLaTeX with its default Latin Modern text font, reports no missing character.
\begin{equation}\label{eq: reg}
    \begin{cases}
        \p_t f_\ve + v\cdot \nabla_x f_\ve = \left(\nu_{f_\ve} +  \ve\right) \nabla_v\cdot \left( \calT_{f_\ve}^{(\ve)} \nabla_v f_\ve +  (\renorm{v}{1}{\ve} - u_{f_\ve}^{(\ve)}) f_\ve \right), \\
        f_\ve|_{t=0} = f_\ve^0 \in C_c(\overline{\Omega}\times\R^d;\R_+), \\
        \gamma f_{\ve}|_{\Sigma_-^T} = g_\ve \in C_c([0,T]\times \p\Omega \times \R^d; \R_+),
    \end{cases}
\end{equation}

\begin{lemma} \label{lem: ve weak sol}
    There exists a weak solution $(f_\ve, \gamma f_\ve)$ to the equation \eqref{eq: reg}, in the sense that
    \begin{align*}
        (f_\ve, \gamma f_\ve) \in L^2((0,T)\times \Omega; H^1_v(\R^d)) \times L^2_{\rm loc}(\Sigma^T, |n(x)\cdot v| \dt\tnd\sigma\dv)
    \end{align*}
    satisfies for each $t\in (0,T]$ and $\varphi\in C_c^\infty([0,t]\times\overline\Omega\times\R^d)$:
    \begin{equation*}
    \begin{split}
    &\int_{\calO} f_\ve(t) \varphi(t,x,v) \,\dx\dv - \int_{\calO} f_\ve^0 \varphi(0,x,v)\,\dx\dv + \int_{\Sigma_+^t} \gamma f_\ve \, \varphi \, (n(x)\cdot v)_+ \ds\tnd\sigma\dv\\
    &\quad = - \int_{\Sigma_-^t} g_\ve \, \varphi \, (n(x)\cdot v)_- \ds\tnd\sigma\dv +\int_{(0,t)\times \calO} f_\ve (\p_s \varphi + v\cdot \nabla_x \varphi) \,\ds\dx\dv\\
    &\qquad + \int_{(0,t)\times \calO} (\nu_{f_\ve} + \ve) \lt(\Delta_v \varphi \calT_{f_\ve}^{(\ve)} f_\ve - \nabla_v\varphi \cdot \lt(\renorm{v}{1}{\ve}-u_{f_\ve}^{(\ve)} \rt) f_\ve \rt) \,\ds\dx\dv.
    \end{split}
    \end{equation*}
    Furthermore, the solution is renormalized, and for each $k\in \bbN$ it satisfies for some constant $C_{T,d,k,\ve}>0$
    \begin{align}\label{eq: ve decay}
        f_\ve(t,x,v) \le \frac{C_{T,d,k,\ve}}{(1+|v|)^k}.
    \end{align}
    In particular,
    \begin{equation}\label{eq: ve memberships}
    \begin{split}
        f_\ve \in L^\infty(0,T; L^1\cap L^\infty(\calO)),  \quad \gamma f_\ve \in L^1 \cap L^\infty (\Sigma^T, |n(x)\cdot v| \dt \tnd\sigma \dv).
        \end{split}
    \end{equation}
\end{lemma}

\begin{lemma} \label{lem: subcubic}
    Let $(f_\ve,\gamma f_\ve)$ denote the solution to \eqref{eq: reg} as constructed through Lemma \ref{lem: ve weak sol}. Then for any $\varphi\in C^0(\R^d)$ with subcubic growth, the averages $\lt\{\intr f_\ve \varphi(v)\,\dv\rt\}_\ve$ are compact in $L^1((0,T)\times\Omega)$.
\end{lemma}

    \begin{equation} \label{eq: MAC convergence}
        \rho_{f_\ve} \to \rho_f, \quad j_{f_\ve} \to j_f, \quad  \intr |v|^2 f_\ve \,\dv \to \intr |v|^2 f \,\dv.
    \end{equation}

    \begin{align}\label{eq: V convergence}
        V_{f_\ve} \to V_f \quad \text{a.e. and in}\quad L^1((0,T)\times\Omega).
    \end{align}

    \begin{align}\label{eq: NU convergence}
        \nu_{f_\ve} + \ve = \nu(\rho_{f_\ve}, j_{f_\ve}, V_{f_\ve}) \to \nu(\rho_f, j_f, V_f) = \nu_f  \quad \text{a.e. in }(0,T)\times \Omega.
    \end{align}

\begin{lemma}
For each $\varphi\in C_c^\infty([0,t]\times \overline{\Omega}\times \R^d)$ we have
\begin{align}\label{eq: PASS weak form 1}
    \int_{(0,t)\times \calO} (\nu_{f_\ve} + \ve) \Delta_v \varphi \calT_{f_\ve}^{(\ve)} f_\ve \,\ds\dx\dv  \to \int_{(0,t)\times \calO} \nu_{f} \Delta_v \varphi \calT_f f \,\ds\dx\dv.
\end{align}
\end{lemma}

\begin{proof}
Resorting to a density argument, we may assume without loss of generality that the test function $\varphi\in C_c^\infty([0,t]\times\overline\Omega\times\R^d)$ takes the form
\[
    \varphi(s,x,v) = \psi(s,x) \eta(v), \quad \psi\in C^\infty([0,t]\times \overline{\Omega}), \quad \eta\in C_c^\infty(\R^d).
\]
Then, with the notations
\[
 \rho_{\Delta\eta}^\ve(s,x) := \intr f_\ve(s,x,v) \Delta\eta(v) \,\dv, \quad \rho_{\Delta\eta}(s,x) := \intr f(s,x,v) \Delta\eta(v)\,\dv.
\]
Note that \eqref{eq: PASS weak form 1} is equivalent to showing that
\begin{align} \label{eq: PASS weak form 1, EQUIV}
    \int_{(0,t)\times \Omega} \psi(s,x) \lt[ (\nu_{f_\ve} + \ve) \calT_{f_\ve}^{(\ve)} \rho_{\Delta\eta}^\ve - \nu_f \calT_f \rho_{\Delta\eta} \rt] \,\ds\dx\dv \to 0.
\end{align}
Hence we only need show that $(\nu_{f_\ve} + \ve) \calT_{f_\ve}^{(\ve)} \rho_{\Delta\eta}^\ve \weakto \nu_f \calT_f \rho_{\Delta\eta}$ weakly in $L^1([0,T]\times \Omega)$. We will prove a stronger statement and show that
\begin{align}\label{eq: stronger}
     \lt[ (\nu_{f_\ve} + \ve) \calT_{f_\ve}^{(\ve)} \rho_{\Delta\eta}^\ve - \nu_f \calT_f \rho_{\Delta\eta} \rt] \to 0 \quad \text{in} \quad L^1((0,T)\times\Omega).
\end{align}
By the averaging effect of Lemma \ref{lem: subcubic}, up to some subsequence it holds that
\begin{align} \label{eq: RHO ETA}
    \rho_{\Delta\eta}^\ve \xrightarrow[\ve\downarrow 0]{} \rho_{\Delta\eta} := \intr f\,\Delta\eta(v)\,\dv \quad \text{a.e. and in} \quad L^1((0,T)\times\Omega).
\end{align}
Taking the product with the (regularized) temperature, we can then observe the following bound, due to Young's inequality
\begin{align}\label{eq: eta YOUNG}
    \calT_{f_\ve}^{(\ve)} |\rho_{\Delta\eta}^\ve| \le C \|\Delta\eta\|_{L^\infty} \lt(1 + \calT_{f_\ve}^{(\ve)}\rt)  \rho_{f_\ve} = C(\|\Delta\eta\|_{L^\infty}) (\rho_{f_\ve} + V_{f_\ve}),
\end{align}
the right-hand side of which is strongly convergent in $L^1((0,T)\times\Omega)$ by \eqref{eq: MAC convergence}, \eqref{eq: V convergence}. As a result, the family $\lt\{\calT_{f_\ve}^{(\ve)} \rho_{\Delta\eta}^\ve \rt\}_\ve$ is equiintegrable in $(0,T)\times\Omega$. Since \eqref{eq: RHO ETA} implies
\begin{align*}
    \calT_{f_\ve}^{(\ve)} \rho_{\Delta\eta}^\ve \xrightarrow[\ve\downarrow 0]{} \calT_f \rho_{\Delta\eta} \quad \text{a.e. in}\quad D_+ := \{(t,x)\in (0,T)\times\Omega: \rho_f(t,x)>0\},
\end{align*}
we deduce from the Vitali convergence theorem that
\begin{align}\label{eq: T ETA D+}
    \calT_{f_\ve}^{(\ve)} \rho_{\Delta\eta}^\ve \to \calT_{f} \rho_{\Delta\eta} \quad \text{in} \quad L^1(D_+).
\end{align}
On the other hand, on the set $D_0 := [(0,T)\times \Omega]\setminus D_+$, we have $f(t,x,v) = 0$ a.e. $v\in \R^d$, which gives that $\rho_{\Delta\eta} = \calT_f = 0$ in $D_0$. Therefore, 
\begin{equation} \label{eq: T ETA D0}
\begin{split}
    \lim_{\ve\downarrow 0} \int_{D_0} \lt| \calT_{f_\ve}^{(\ve)} \rho_{\Delta\eta}^\ve - \calT_f \rho_{\Delta\eta} \rt| \,\dt\dx &= \lim_{\ve\downarrow 0}\int_{D_0} \calT_{f_\ve}^{(\ve)} \lt|\rho_{\Delta\eta}^\ve \rt| \,\dt\dx \\
    &\le C(\|\Delta\eta\|_{L^\infty}) \lim_{\ve\downarrow 0} \int_{D_0} (\rho_{f_\ve} + V_{f_\ve}) \,\dt\dx\\
    &= C \int_{D_0} (\rho_f + V_f) \,\dt\dx\\
    &= 0,
\end{split}
\end{equation}
the inequality due to \eqref{eq: eta YOUNG}. In other words this proves $\calT_{f_\ve}^{(\ve)} \rho_{\Delta\eta}^\ve \to \calT_f \rho_{\Delta\eta}$ in $L^1(D_0)$. Combining this with \eqref{eq: T ETA D+}, we find
\begin{align}\label{eq: PASS prod T eta}
    \calT_{f_\ve}^{(\ve)} \rho_{\Delta\eta}^\ve \to \calT_f \rho_{\Delta\eta} \quad \text{in} \quad L^1((0,T)\times\Omega).
\end{align}
To complete the proof, we recall that $(\nu_{f_\ve} + \ve) \to \nu_f$ pointwise a.e. in $(0,T)\times \Omega$ (see \eqref{eq: NU convergence}), and $\|\nu_{f_\ve} + \ve\|_{L^\infty} \le \|\nu\|_{L^\infty} + 1$. By Egorov's theorem we can extract a set $E\subset (0,T)\times\Omega$ with arbitrarily small Lebesgue measure for which $\nu_{f_\ve} \to \nu_f$ uniformly outside $E$. Then, writing the difference in \eqref{eq: PASS weak form 1, EQUIV} as
\begin{equation} \label{eq: j1 j2}
\begin{split}
     \lt|(\nu_{f_\ve} + \ve) \calT_{f_\ve}^{(\ve)} \rho_{\Delta\eta}^\ve - \nu_f \calT_f \rho_{\Delta\eta} \rt|  
    &\le \lt| (\nu_{f_\ve} + \ve - \nu_f)  \calT_{f_\ve}^{(\ve)} \rho_{\Delta\eta}^\ve \rt| + \lt|\nu_f \lt( \calT_{f_\ve}^{(\ve)} \rho_{\Delta\eta}^\ve - \calT_f \rho_{\Delta\eta} \rt) \rt|\\
    &=: J_1 + J_2,
\end{split}
\end{equation}
we see that $J_1\to 0$ in $L^1([(0,T)\times\Omega]\setminus E)$ and $J_2\to 0$ in $L^1((0,T)\times\Omega)$, thanks to \eqref{eq: PASS prod T eta}. Taking the Lebesgue measure of $E$ to zero, we then see that $\|J_1\|_{L^1(E)}\to 0$ since $\nu$ is bounded and the family $\lt\{\calT_{f_\ve}^{(\ve)} \rho_{\Delta\eta}^\ve\rt\}$ is equiintegrable. It follows that
\begin{align*}
    (\nu_{f_\ve} + \ve) \calT_{f_\ve}^{(\ve)} \rho_{\Delta\eta}^\ve \to \nu_f \calT_f \rho_{\Delta\eta} \quad \text{in}\quad L^1((0,T)\times \Omega),
\end{align*}
and this proves \eqref{eq: stronger}. Clearly, the latter then implies \eqref{eq: PASS weak form 1, EQUIV}.
\end{proof}

\end{document}
